% Worked example. Packages: booktabs, amsmath. Source: docs/EXAMPLES.txt, produced by
%   python3 scripts/worked-example.py results/xp-main/soybean/taste1/m3-gain/seed3.trace.jsonl -it 0 1 2 --liked "29,37;33,45"
% Iteration t here is iteration t-1 in the trace.

\subsection{Worked Example}
\label{subsec:worked_example}

We follow the first three iterations of one run: soybean, a combination user who likes the item pairs
$\{29, 37\}$ (\texttt{stem\_cankers=absent}, \texttt{external\_decay=absent}) and $\{33, 45\}$
(\texttt{canker\_lesion=dna}, \texttt{sclerotia=absent}), query size $k = 10$, $m = 3$ selected with
\textsc{PartitionGain}, and the published folding scheme. Ranks go from $0$ (preferred) to $9$, so
$\mu = 4.5$ and $\mathrm{TSS} = \sum_{r=0}^{9} (r - 4.5)^2 = 82.5$.

\paragraph{Iteration 1: query and ranking.}
Table~\ref{tab:ex_query1} shows the ten sampled patterns in the order the user ranked them. Since all
weights start at $1$, the model score $w \cdot x$ of a pattern is its length. The user puts the three patterns
containing $\{33, 45\}$ first (score $1$ plus the small surprisingness tie-break); the others are ordered by
surprisingness alone.

\begin{table}[t]
\centering
\caption{Iteration 1: the query, ranked by the user. \#: position in the sampled order.}
\label{tab:ex_query1}
\small
\begin{tabular}{rrlrr}
\toprule
Rank & \# & Pattern & $w \cdot x$ & User score \\
\midrule
0 & 3 & $\{8,11,27,29,\mathbf{33},40,42,\mathbf{45}\}$ & 8 & 1.0008 \\
1 & 6 & $\{19,\mathbf{33},42,\mathbf{45},47\}$          & 5 & 1.0003 \\
2 & 7 & $\{7,13,18,27,\mathbf{33},\mathbf{45}\}$        & 6 & 1.0002 \\
3 & 5 & $\{11,26,28,32,40\}$                            & 5 & 0.0012 \\
4 & 4 & $\{11,13,28,35,40,42,45\}$                      & 7 & 0.0007 \\
5 & 1 & $\{4,7,23,26,40,45\}$                           & 6 & 0.0004 \\
6 & 9 & $\{10,20,35,40,42\}$                            & 5 & 0.0003 \\
7 & 8 & $\{11,13,23,26,34,40,47\}$                      & 7 & 0.0003 \\
8 & 2 & $\{14,18,27,37,40,42,47\}$                      & 7 & 0.0003 \\
9 & 0 & $\{7,11,23,32,37,40,45\}$                       & 7 & 0.0000 \\
\bottomrule
\end{tabular}
\end{table}

\paragraph{Iteration 1: ICV.}
Among the sub-itemsets of the query, $39$ distinct covers have a positive ICV. The highest is that of the liked
pair $Y = \{33, 45\}$, contained in the patterns of ranks $0, 1, 2$:
\begin{equation*}
  \bar{r}_{C_Y} = \tfrac{0+1+2}{3} = 1, \qquad
  \bar{r}_{\overline{C}_Y} = \tfrac{3+4+\dots+9}{7} = 6, \qquad
  \mathrm{ICV}(Y) = 3\,(4.5 - 1)^2 + 7\,(4.5 - 6)^2 = 36.75 + 15.75 = 52.5 .
\end{equation*}
The next candidates are $\{33, 42, 45\}$ (ranks $0, 1$; ICV $40.0$), $\{37, 40\}$ (ranks $8, 9$; ICV $40.0$) and
$\{42, 45\}$ (ranks $0, 1, 4$; ICV $34.4$).

\paragraph{Iteration 1: partition gain.}
\textsc{PartitionGain} starts from $\{33, 45\}$, so $\mathrm{PICV} = 52.5$, and then adds the candidate that
most refines the partition of the ranks (Table~\ref{tab:ex_gain1}). $\{33, 42, 45\}$ has the second-highest ICV
but adds nothing: its cover $\{0, 1\}$ lies inside that of $\{33, 45\}$ and only splits the preferred group
further, whereas $\{37, 40\}$ separates the two least preferred patterns. In step 2, $\{37, 40\}$ and
$\{11, 28, 40\}$ tie with a gain of $17.5$; the tie goes to the higher ICV.

\begin{table}[t]
\centering
\caption{Iteration 1: greedy partition gain (best candidates per step; selected in bold).}
\label{tab:ex_gain1}
\small
\begin{tabular}{clcccc}
\toprule
Step & Candidate & Ranks covered & ICV & PICV & Gain \\
\midrule
1 & $\mathbf{\{33,45\}}$     & 0, 1, 2 & 52.5 & 52.5 & -- \\
\midrule
2 & $\mathbf{\{37,40\}}$     & 8, 9    & 40.0 & 70.0 & +17.5 \\
  & $\{11,28,40\}$           & 3, 4    & 2.5  & 70.0 & +17.5 \\
  & $\{11,23,40\}$           & 7, 9    &      & 63.7 & +11.2 \\
\midrule
3 & $\mathbf{\{11,28,40\}}$  & 3, 4    & 2.5  & 77.5 & +7.5 \\
  & $\{11,23,40\}$           & 7, 9    &      & 75.5 & +5.5 \\
  & $\{40,47\}$              & 7, 8    &      & 75.5 & +5.5 \\
\bottomrule
\end{tabular}
\end{table}

The selected set partitions the query into four groups (Table~\ref{tab:ex_partition1}), which explain
$77.5$ of the $82.5$ total: the preferred patterns, the two least preferred ones, ranks $3$--$4$, and the
remaining ranks $5$--$7$.

\begin{table}[t]
\centering
\caption{Iteration 1: final partition. Signature: which of $\{33,45\}$, $\{37,40\}$, $\{11,28,40\}$ the
pattern contains.}
\label{tab:ex_partition1}
\small
\begin{tabular}{cccc}
\toprule
Signature & Ranks & Mean rank & $n_g (\mu - \bar{r}_g)^2$ \\
\midrule
100 & 0, 1, 2 & 1.0 & 36.75 \\
010 & 8, 9    & 8.5 & 32.00 \\
001 & 3, 4    & 3.5 & 2.00 \\
000 & 5, 6, 7 & 6.0 & 6.75 \\
\midrule
    &         &     & $\mathrm{PICV} = 77.5$ \\
\bottomrule
\end{tabular}
\end{table}

\paragraph{Iterations 2 and 3.}
Table~\ref{tab:ex_iter23} summarises the next two iterations. The model scores are no longer equal to the
pattern lengths, since the learned weights now drive the sampler. In both queries the user's other liked pair,
$\{29, 37\}$, is the highest-ICV sub-pattern ($52.5$ and $40.0$). The other selected sub-patterns
are chosen for the groups they separate, not for their own ICV. In iteration 2, $\{45\}$ (ICV $9.6$) is added
before candidates with a far higher ICV, such as $\{26, 29, 37, 45\}$ (ICV $40.0$), which only splits the
preferred group again. Its gain of $19.0$ comes from separating ranks $5$--$9$ from ranks $3$--$4$. In step 3,
$\{40\}$ and $\{11\}$ tie: their covers are complementary, so they induce the same partition and have the same
ICV and length. The program keeps the one found first.

\begin{table}[t]
\centering
\caption{Iterations 2 and 3: ranked queries (liked pair in bold) and selected sub-patterns.}
\label{tab:ex_iter23}
\small
\begin{tabular}{rlrr@{\hskip 2em}rlrr}
\toprule
\multicolumn{4}{c}{Iteration 2} & \multicolumn{4}{c}{Iteration 3} \\
\cmidrule(r){1-4} \cmidrule(l){5-8}
Rank & Pattern & $w\cdot x$ & Score & Rank & Pattern & $w\cdot x$ & Score \\
\midrule
0 & $\{8,21,26,27,\mathbf{29},\mathbf{37},40,45\}$ & 6.99 & 1.0003 & 0 & $\{8,18,26,27,\mathbf{29},\mathbf{37}\}$ & 5.42 & 1.0004 \\
1 & $\{7,11,14,26,\mathbf{29},\mathbf{37},45\}$    & 5.40 & 1.0002 & 1 & $\{4,23,\mathbf{29},\mathbf{37}\}$       & 2.89 & 1.0004 \\
2 & $\{10,13,\mathbf{29},\mathbf{37},40,42,47\}$   & 5.63 & 1.0002 & 2 & $\{20,24,26,35\}$                      & 3.66 & 0.0012 \\
3 & $\{8,40,42,48\}$                               & 4.03 & 0.0007 & 3 & $\{11,23,33,37,40\}$                   & 3.76 & 0.0010 \\
4 & $\{11,23,32\}$                                 & 1.72 & 0.0005 & 4 & $\{23,26,32,35,42,47\}$                & 4.84 & 0.0005 \\
5 & $\{19,26,38,40,42,45,47\}$                     & 6.97 & 0.0003 & 5 & $\{11,14,23,26,28,38,40,47\}$          & 5.61 & 0.0004 \\
6 & $\{3,8,24,26,28,40,45\}$                       & 7.25 & 0.0003 & 6 & $\{7,13,22,26,33,37,40,42,47\}$        & 8.32 & 0.0004 \\
7 & $\{7,16,20,23,40,42,45\}$                      & 5.78 & 0.0002 & 7 & $\{7,33,42\}$                          & 4.36 & 0.0003 \\
8 & $\{5,11,23,26,28,45\}$                         & 4.84 & 0.0000 & 8 & $\{14,17,23,33,40\}$                   & 4.41 & 0.0001 \\
9 & $\{7,11,13,35,45\}$                            & 4.21 & 0.0000 & 9 & $\{11,13,18,26,27,37\}$                & 3.94 & 0.0001 \\
\midrule
\multicolumn{4}{l}{Selected (ICV; PICV after adding):} & \multicolumn{4}{l}{Selected (ICV; PICV after adding):} \\
\multicolumn{4}{l}{$\{29,37\}$ ($52.5$; $52.5$), $\{45\}$ ($9.6$; $71.5$), $\{11\}$ ($6.7$; $80.0$)} &
\multicolumn{4}{l}{$\{29,37\}$ ($40.0$; $40.0$), $\{26,35\}$ ($5.6$; $56.7$), $\{11,23,40\}$ ($0.6$; $73.0$)} \\
\bottomrule
\end{tabular}
\end{table}

\paragraph{What happens next.}
With folding, the weights learned for these sub-patterns are transferred to their items and the sub-pattern
features are removed, so the preference for $29$ and $37$ \emph{together} is lost and the pair has to be
rediscovered in later queries. With \texttt{keep}, $\{29, 37\}$ and $\{33, 45\}$ stay in the model as features.
In the keep run with the same seed, the weight of $\{29, 37\}$ grows from $0.26$ after iteration 2 to $1.16$
after iteration 20, the largest of the $47$ kept sub-patterns, and that of $\{33, 45\}$ from $0$ to $0.72$
(Section~\ref{subsec:keep_results}).
